\documentclass[aspectratio=169,11pt]{beamer}
\usepackage[T1]{fontenc}
\usepackage[utf8]{inputenc}
\usepackage{lmodern,microtype,booktabs,tabularx,array,ragged2e,xcolor,hyperref}

\definecolor{StudentNavy}{HTML}{102A43}
\definecolor{StudentBlue}{HTML}{1769E0}
\definecolor{StudentCyan}{HTML}{14A6B8}
\definecolor{StudentGold}{HTML}{E8A628}
\definecolor{StudentInk}{HTML}{243B53}
\definecolor{StudentMist}{HTML}{F2F7FC}
\definecolor{StudentSoft}{HTML}{D9E8F5}
\definecolor{StudentGray}{HTML}{627D98}
\definecolor{StudentWhite}{HTML}{FFFFFF}
\definecolor{StudentGreen}{HTML}{2B8A6E}

\hypersetup{colorlinks=true,urlcolor=StudentBlue,linkcolor=StudentBlue}
\setbeamercolor{normal text}{fg=StudentInk,bg=StudentWhite}
\setbeamercolor{structure}{fg=StudentBlue}
\setbeamercolor{frametitle}{fg=StudentNavy,bg=StudentWhite}
\setbeamercolor{title}{fg=StudentWhite}
\setbeamercolor{subtitle}{fg=StudentSoft}
\setbeamercolor{author}{fg=StudentWhite}
\setbeamercolor{date}{fg=StudentSoft}
\setbeamercolor{block title}{fg=StudentWhite,bg=StudentBlue}
\setbeamercolor{block body}{fg=StudentInk,bg=StudentMist}
\setbeamerfont{title}{series=\bfseries,size=\fontsize{30}{35}\selectfont}
\setbeamerfont{subtitle}{size=\fontsize{15}{19}\selectfont}
\setbeamerfont{author}{series=\bfseries,size=\fontsize{15}{18}\selectfont}
\setbeamerfont{date}{size=\fontsize{11.5}{14}\selectfont}
\setbeamerfont{frametitle}{series=\bfseries,size=\fontsize{19.5}{23}\selectfont}
\setbeamerfont{normal text}{size=\fontsize{12.2}{15}\selectfont}
\setbeamerfont{block title}{series=\bfseries,size=\fontsize{13}{16}\selectfont}
\setbeamerfont{block body}{size=\fontsize{12}{15.5}\selectfont}
\setbeamerfont{footline}{size=\fontsize{8}{10}\selectfont}
\setbeamertemplate{navigation symbols}{}
\setbeamertemplate{items}[circle]
\setbeamertemplate{itemize item}{\color{StudentCyan}\raisebox{0.12ex}{\scriptsize$\blacktriangleright$}}
\setbeamertemplate{itemize subitem}{\color{StudentBlue}\scriptsize$\bullet$}
\setlength{\leftmargini}{1.2em}
\setlength{\leftmarginii}{1.1em}

\setbeamertemplate{frametitle}{%
  \nointerlineskip
  \begin{beamercolorbox}[wd=\paperwidth,leftskip=0.7cm,rightskip=0.7cm,sep=0.17cm]{frametitle}
    \rightskip=0pt plus 1fil\relax\usebeamerfont{frametitle}\insertframetitle\par
    \color{StudentCyan}\rule{\dimexpr\paperwidth-1.8cm\relax}{0.8pt}
  \end{beamercolorbox}
}
\setbeamertemplate{footline}{%
  \leavevmode\hbox{%
    \begin{beamercolorbox}[wd=.67\paperwidth,ht=2.7ex,dp=1.2ex,leftskip=.7cm]{author in head/foot}
      \color{StudentGray}\usebeamerfont{footline}Sosialisasi Orang Tua \textbar\ SMAN 1 Muntilan
    \end{beamercolorbox}
    \begin{beamercolorbox}[wd=.31\paperwidth,ht=2.7ex,dp=1.2ex,center]{date in head/foot}
      \color{StudentBlue}\usebeamerfont{footline}Bab 4 \textbar\ \insertframenumber/\inserttotalframenumber
    \end{beamercolorbox}
  }
}
\newcommand{\kicker}[1]{{\color{StudentCyan}\bfseries\fontsize{10.5}{13}\selectfont\MakeUppercase{#1}}\par\vspace{0.11cm}}
\newcommand{\claim}[1]{{\color{StudentNavy}\bfseries\fontsize{16}{19}\selectfont #1}\par}
\newcommand{\metric}[2]{{\color{StudentBlue}\bfseries\fontsize{27}{30}\selectfont #1}\par\vspace{0.04cm}{\color{StudentGray}\fontsize{11.5}{14.5}\selectfont #2}}

\title{Pelaksanaan Pembimbingan UTBK 2026}
\subtitle{Bab 4 \textbar\ Dari materi, latihan, dan tryout menuju tindak lanjut}
\author{Balya Rochmadi}
\date{SMAN 1 Muntilan bersama StudentPro \textbar\ Tahun Pelajaran 2025/2026}

\begin{document}

\begingroup
\setbeamercolor{background canvas}{bg=StudentNavy}
\begin{frame}[plain]
  \vspace{0.65cm}
  \kicker{SMAN 1 Muntilan \textbar\ StudentPro}
  \vspace{0.08cm}
  {\usebeamerfont{title}\color{StudentWhite}\inserttitle\par}
  \vspace{0.38cm}
  {\usebeamerfont{subtitle}\color{StudentSoft}\insertsubtitle\par}
  \vfill
  {\color{StudentGold}\rule{2.3cm}{2.5pt}}\par
  \vspace{0.20cm}
  {\usebeamerfont{author}\color{StudentWhite}\insertauthor}\par
  {\usebeamerfont{date}\color{StudentSoft}\insertdate}
  \vspace{0.42cm}
\end{frame}
\endgroup

\begin{frame}{Bab 4 mencatat proses yang benar-benar dilaksanakan}
  \kicker{Batas Pembahasan}
  \begin{columns}[T,onlytextwidth]
    \begin{column}{0.31\textwidth}
      \begin{block}{Masuk Bab 4}
        Peserta, alur kegiatan, materi, latihan, tryout, analisis, pembahasan, konsultasi, dan monitoring.
      \end{block}
    \end{column}
    \begin{column}{0.31\textwidth}
      \begin{block}{Masuk Bab 5}
        Hasil akhir, perubahan skor, kekuatan, kelemahan, kesenjangan, dan dampak program.
      \end{block}
    \end{column}
    \begin{column}{0.31\textwidth}
      \begin{block}{Prinsip pelaporan}
        Angka dicantumkan apabila terdapat catatan atau dokumen internal pendukung.
      \end{block}
    \end{column}
  \end{columns}
  \vspace{0.24cm}
  \centering{\color{StudentNavy}\bfseries\fontsize{14}{17}\selectfont Bab ini menjelaskan apa yang dikerjakan sebelum menilai hasilnya.}
\end{frame}

\begin{frame}[shrink=6]{Program menjangkau peserta dengan kebutuhan beragam}
  \kicker{Sasaran Tahun Pelajaran 2025/2026}
  \begin{columns}[T,onlytextwidth]
    \begin{column}{0.30\textwidth}
      \begin{block}{Siswa kelas XII}
        Peserta utama yang memerlukan persiapan akademik, strategi ujian, dan pilihan program studi.
      \end{block}
    \end{column}
    \begin{column}{0.30\textwidth}
      \begin{block}{Alumni}
        Peserta yang kembali menyiapkan UTBK dengan pengalaman ujian dan kebutuhan remedial berbeda.
      \end{block}
    \end{column}
    \begin{column}{0.30\textwidth}
      \begin{block}{Calon peserta}
        Peserta layanan yang membutuhkan pemetaan sebelum mengikuti rangkaian pembimbingan.
      \end{block}
    \end{column}
  \end{columns}
  \vspace{0.26cm}
  \claim{Satu program digunakan bersama, tetapi tindak lanjutnya tidak disamaratakan.}
\end{frame}

\begin{frame}[shrink=3]{Pelaksanaan dibangun sebagai satu siklus belajar}
  \kicker{Arsitektur Program}
  {\fontsize{10.2}{12.7}\selectfont
  \renewcommand{\arraystretch}{0.99}
  \begin{tabularx}{\textwidth}{>{\bfseries\color{StudentBlue}}p{0.18\textwidth} >{\RaggedRight\arraybackslash}X >{\RaggedRight\arraybackslash}p{0.30\textwidth}}
    \toprule
    \textbf{Tahap} & \textbf{Kegiatan utama} & \textbf{Keluaran} \\
    \midrule
    Pendataan & Menghimpun peserta dan kebutuhan awal & Daftar peserta dan status layanan \\
    Diagnostik & Mengukur kemampuan awal dan kebiasaan belajar & Peta awal per siswa \\
    Pembelajaran & Materi, bank soal, latihan, dan pembahasan & Catatan penguasaan dan kesalahan \\
    Simulasi & Tryout berkala dengan batas waktu & Skor, waktu, dan pola pengerjaan \\
    Tindak lanjut & Konsultasi, remedial, strategi, dan monitoring & Rencana belajar berikutnya \\
    \bottomrule
  \end{tabularx}}
  \vspace{0.10cm}
  \centering{\color{StudentNavy}\bfseries\fontsize{10.8}{13.4}\selectfont Siklus diulang agar data latihan menghasilkan tindakan.}
\end{frame}

\begin{frame}[shrink=3]{Pendataan dan diagnostik menentukan titik awal}
  \kicker{Tahap 1 -- Pemetaan}
  \begin{columns}[T,onlytextwidth]
    \begin{column}{0.46\textwidth}
      {\color{StudentBlue}\bfseries Data yang dihimpun}\par
      \begin{itemize}
        \item Status siswa dan rencana mengikuti UTBK.
        \item Riwayat belajar serta kebutuhan pendampingan.
        \item Hasil diagnostik awal.
        \item Arah program studi dan perguruan tinggi.
      \end{itemize}
    \end{column}
    \begin{column}{0.48\textwidth}
      {\color{StudentBlue}\bfseries Keputusan awal}\par
      \begin{itemize}
        \item Urutan subtes yang perlu diperbaiki.
        \item Materi dasar yang harus diulang.
        \item Jadwal latihan dan tryout.
        \item Kebutuhan konsultasi individual.
      \end{itemize}
    \end{column}
  \end{columns}
  \vspace{0.18cm}
  \begin{beamercolorbox}[rounded=true,sep=0.18cm,wd=\textwidth]{block body}
    \centering\bfseries Diagnostik digunakan untuk mengarahkan belajar, bukan memberi label kemampuan tetap.
  \end{beamercolorbox}
\end{frame}

\begin{frame}{Pembimbingan materi mengikuti kebutuhan tujuh subtes}
  \kicker{Tahap 2 -- Penguatan Akademik}
  \begin{columns}[T,onlytextwidth]
    \begin{column}{0.46\textwidth}
      {\color{StudentBlue}\bfseries Tes Potensi Skolastik}\par
      \begin{enumerate}
        \item Penalaran Umum
        \item Pengetahuan dan Pemahaman Umum
        \item Pemahaman Bacaan dan Menulis
        \item Pengetahuan Kuantitatif
      \end{enumerate}
    \end{column}
    \begin{column}{0.48\textwidth}
      {\color{StudentBlue}\bfseries Tes Literasi}\par
      \begin{enumerate}
        \item Literasi dalam Bahasa Indonesia
        \item Literasi dalam Bahasa Inggris
        \item Penalaran Matematika
      \end{enumerate}
    \end{column}
  \end{columns}
  \vspace{0.20cm}
  \begin{beamercolorbox}[rounded=true,sep=0.18cm,wd=\textwidth]{block body}
    \centering\bfseries Materi menjembatani konsep dasar menuju penerapan dalam soal.
  \end{beamercolorbox}
\end{frame}

\begin{frame}[shrink=5]{Materi dan bank soal didistribusikan sebagai satu rangkaian}
  \kicker{Bukan Hanya Menambah Jumlah Soal}
  {\fontsize{10.4}{13}\selectfont
  \renewcommand{\arraystretch}{1.02}
  \begin{tabularx}{\textwidth}{>{\bfseries\color{StudentBlue}}p{0.22\textwidth} >{\RaggedRight\arraybackslash}X >{\RaggedRight\arraybackslash}p{0.31\textwidth}}
    \toprule
    \textbf{Komponen} & \textbf{Fungsi} & \textbf{Bukti belajar} \\
    \midrule
    Materi ringkas & Menata konsep dan strategi pokok & Siswa mengetahui apa yang harus dikuasai \\
    Contoh soal & Memperlihatkan penerapan langkah berpikir & Siswa dapat mengikuti alasan penyelesaian \\
    Bank soal & Memberi variasi indikator dan tingkat kesulitan & Akurasi dan pola kesalahan mulai terlihat \\
    Pembahasan & Mengoreksi alasan, bukan hanya jawaban akhir & Siswa dapat menjelaskan perbaikannya \\
    \bottomrule
  \end{tabularx}}
  \vspace{0.12cm}
  \centering{\color{StudentNavy}\bfseries\fontsize{10.8}{13.5}\selectfont Latihan bernilai ketika kesalahan dapat dikenali dan diperbaiki.}
\end{frame}

\begin{frame}{Tryout berkala menguji lebih dari penguasaan materi}
  \kicker{Tahap 3 -- Simulasi}
  \begin{columns}[T,onlytextwidth]
    \begin{column}{0.30\textwidth}
      \begin{block}{Akurasi}
        Seberapa banyak jawaban benar dan materi apa yang masih lemah?
      \end{block}
    \end{column}
    \begin{column}{0.30\textwidth}
      \begin{block}{Waktu}
        Apakah siswa terlalu lama, terburu-buru, atau tidak selesai?
      \end{block}
    \end{column}
    \begin{column}{0.30\textwidth}
      \begin{block}{Stamina}
        Apakah kualitas keputusan bertahan sampai rangkaian tes selesai?
      \end{block}
    \end{column}
  \end{columns}
  \vspace{0.26cm}
  \claim{Tryout adalah alat diagnosis, bukan sekadar perlombaan peringkat.}
\end{frame}

\begin{frame}[shrink=12]{Scoring diteruskan menjadi analisis yang dapat ditindaklanjuti}
  \kicker{Setelah Tryout}
  {\fontsize{10.1}{12.6}\selectfont
  \renewcommand{\arraystretch}{0.98}
  \begin{tabularx}{\textwidth}{>{\bfseries\color{StudentBlue}}p{0.20\textwidth} >{\RaggedRight\arraybackslash}p{0.31\textwidth} >{\RaggedRight\arraybackslash}X}
    \toprule
    \textbf{Data} & \textbf{Pertanyaan analisis} & \textbf{Tindakan} \\
    \midrule
    Skor per subtes & Bagian mana yang paling menahan hasil total? & Menentukan prioritas remedial \\
    Ketepatan materi & Indikator apa yang berulang kali salah? & Mengulang konsep dan pola soal \\
    Waktu pengerjaan & Tidak paham atau terlalu terburu-buru? & Mengubah strategi waktu \\
    Tren beberapa tes & Perbaikan stabil atau hanya terjadi sekali? & Melanjutkan atau mengubah program \\
    \bottomrule
  \end{tabularx}}
  \vspace{0.10cm}
  \centering{\color{StudentNavy}\bfseries\fontsize{10.6}{13.2}\selectfont Nilai berguna setelah menjawab pertanyaan: Apa tindakan berikutnya?}
\end{frame}

\begin{frame}{Pembahasan memperbaiki cara berpikir siswa}
  \kicker{Tahap 4 -- Koreksi dan Remedial}
  \begin{columns}[T,onlytextwidth]
    \begin{column}{0.46\textwidth}
      {\color{StudentBlue}\bfseries Kesalahan yang ditelusuri}\par
      \begin{itemize}
        \item Konsep belum dikuasai.
        \item Salah membaca informasi.
        \item Langkah terlalu panjang.
        \item Strategi tidak efisien.
        \item Kurang teliti atau kehabisan waktu.
      \end{itemize}
    \end{column}
    \begin{column}{0.48\textwidth}
      {\color{StudentBlue}\bfseries Tindak lanjut}\par
      \begin{itemize}
        \item Penjelasan ulang konsep.
        \item Latihan setara pada indikator sama.
        \item Pembiasaan menandai bukti.
        \item Latihan dengan batas waktu.
        \item Pengayaan untuk kemampuan yang kuat.
      \end{itemize}
    \end{column}
  \end{columns}
\end{frame}

\begin{frame}{Konsultasi menjadi bagian penting dari pelaksanaan}
  \kicker{Layanan Individual yang Terdokumentasi}
  \begin{columns}[T,onlytextwidth]
    \begin{column}{0.30\textwidth}
      \metric{124}{siswa tercatat menerima layanan konsultasi akademik}
      \vspace{0.16cm}
      {\color{StudentGreen}\bfseries\fontsize{10.5}{13}\selectfont Sumber: LPJ internal UTBK 2026.}
    \end{column}
    \begin{column}{0.64\textwidth}
      \begin{itemize}
        \item Membaca hasil diagnostik dan tryout.
        \item Menyusun strategi belajar pribadi.
        \item Mengevaluasi hambatan akademik dan kebiasaan belajar.
        \item Membahas program studi dan perguruan tinggi.
        \item Membuat rekap serta rencana tindak lanjut.
      \end{itemize}
    \end{column}
  \end{columns}
  \vspace{0.12cm}
  {\color{StudentGray}\fontsize{9.5}{11.8}\selectfont Angka 124 menyatakan siswa penerima layanan yang terdokumentasi, bukan jumlah seluruh pertemuan konsultasi.}
\end{frame}

\begin{frame}{Pilihan program studi dibahas bersama kesiapan akademik}
  \kicker{Konsultasi Jurusan dan Kampus}
  \begin{columns}[T,onlytextwidth]
    \begin{column}{0.30\textwidth}
      \begin{block}{Minat}
        Bidang yang ingin dipelajari dan arah pekerjaan yang dipertimbangkan.
      \end{block}
    \end{column}
    \begin{column}{0.30\textwidth}
      \begin{block}{Kemampuan}
        Profil subtes, konsistensi latihan, serta kesiapan memperbaiki kelemahan.
      \end{block}
    \end{column}
    \begin{column}{0.30\textwidth}
      \begin{block}{Peluang}
        Daya tampung, persaingan, serta alternatif pilihan yang rasional.
      \end{block}
    \end{column}
  \end{columns}
  \vspace{0.28cm}
  \begin{beamercolorbox}[rounded=true,sep=0.20cm,wd=\textwidth]{block body}
    \centering\bfseries Konsultasi membantu siswa mengambil keputusan dengan lebih sadar.
  \end{beamercolorbox}
\end{frame}

\begin{frame}[shrink=7]{Motivasi dan monitoring menjaga kesinambungan program}
  \kicker{Pendampingan di Antara Pertemuan}
  \begin{columns}[T,onlytextwidth]
    \begin{column}{0.46\textwidth}
      {\color{StudentBlue}\bfseries Monitoring harian}\par
      \begin{itemize}
        \item Kehadiran dan partisipasi.
        \item Penyelesaian materi serta latihan.
        \item Hambatan yang perlu segera ditangani.
        \item Konsistensi jadwal belajar.
      \end{itemize}
    \end{column}
    \begin{column}{0.48\textwidth}
      {\color{StudentBlue}\bfseries Penguatan motivasi}\par
      \begin{itemize}
        \item Menjaga tujuan tetap realistis.
        \item Mengubah kegagalan latihan menjadi bahan belajar.
        \item Mengurangi kecemasan menjelang ujian.
        \item Menjaga tidur, kesehatan, dan ritme belajar.
      \end{itemize}
    \end{column}
  \end{columns}
  \vspace{0.14cm}
  \claim{Konsistensi dibangun melalui perhatian kecil yang dilakukan terus-menerus.}
\end{frame}

\begin{frame}[shrink=8]{StudentPro menghubungkan kegiatan yang sebelumnya terpisah}
  \kicker{Peran LMS dalam Pelaksanaan}
  {\fontsize{10.0}{12.5}\selectfont
  \renewcommand{\arraystretch}{0.98}
  \begin{tabularx}{\textwidth}{>{\bfseries\color{StudentBlue}}p{0.20\textwidth} >{\RaggedRight\arraybackslash}X >{\RaggedRight\arraybackslash}p{0.31\textwidth}}
    \toprule
    \textbf{Fungsi} & \textbf{Kegiatan} & \textbf{Informasi yang tersimpan} \\
    \midrule
    Materi & Distribusi modul per subtes & Akses dan ketuntasan materi \\
    Latihan & Kuis topikal dan bank soal & Jawaban, akurasi, dan indikator \\
    Tryout & Paket subtes dan simulasi penuh & Skor, waktu, serta penyelesaian \\
    Analisis & Rekap hasil dan pola kesalahan & Prioritas remedial/pengayaan \\
    Monitoring & Catatan progres dan konsultasi & Tren dan tindak lanjut siswa \\
    \bottomrule
  \end{tabularx}}
  \vspace{0.08cm}
  \centering{\color{StudentNavy}\bfseries\fontsize{10.4}{13}\selectfont LMS membantu kegiatan belajar meninggalkan jejak data.}
\end{frame}

\begin{frame}[shrink=15]{Alur data membuat pendampingan lebih terarah}
  \kicker{Dari Aktivitas Menjadi Keputusan}
  {\fontsize{10.2}{12.7}\selectfont
  \renewcommand{\arraystretch}{1.00}
  \begin{tabularx}{\textwidth}{>{\bfseries\color{StudentBlue}}p{0.18\textwidth} >{\RaggedRight\arraybackslash}p{0.27\textwidth} >{\RaggedRight\arraybackslash}X}
    \toprule
    \textbf{Masukan} & \textbf{Pengolahan} & \textbf{Keputusan pembimbingan} \\
    \midrule
    Hasil diagnostik & Pemetaan subtes dan indikator & Prioritas materi awal \\
    Aktivitas belajar & Ketuntasan materi dan latihan & Pengingat atau bantuan belajar \\
    Hasil tryout & Skor, waktu, tren, dan kesalahan & Remedial, pengayaan, strategi waktu \\
    Catatan konsultasi & Tujuan, hambatan, dan kesepakatan & Rencana belajar serta pilihan studi \\
    \bottomrule
  \end{tabularx}}
  \vspace{0.14cm}
  \begin{beamercolorbox}[rounded=true,sep=0.17cm,wd=\textwidth]{block body}
    \centering\bfseries Data tidak menggantikan percakapan; data membuat percakapan lebih tepat sasaran.
  \end{beamercolorbox}
\end{frame}

\begin{frame}{Pelaksanaan membutuhkan pembagian peran yang jelas}
  \kicker{Kolaborasi Program 2026}
  {\fontsize{10.3}{12.8}\selectfont
  \renewcommand{\arraystretch}{1.00}
  \begin{tabularx}{\textwidth}{>{\bfseries\color{StudentBlue}}p{0.20\textwidth} >{\RaggedRight\arraybackslash}X}
    \toprule
    \textbf{Pihak} & \textbf{Peran dalam pelaksanaan} \\
    \midrule
    Sekolah & Koordinasi peserta, jadwal, komunikasi, dan dukungan kebijakan. \\
    Tentor & Materi, latihan, pembahasan, analisis, konsultasi, dan tindak lanjut. \\
    Siswa & Mengikuti kegiatan dengan jujur dan menggunakan umpan balik. \\
    Orang tua & Menjaga rutinitas, kesehatan, ketenangan, dan komunikasi. \\
    StudentPro & Menyediakan materi, kuis, tryout, rekaman aktivitas, dan laporan. \\
    \bottomrule
  \end{tabularx}}
\end{frame}

\begin{frame}{Dokumentasi 2026 menjadi dasar perbaikan program}
  \kicker{Bukti Pelaksanaan yang Digunakan}
  \begin{columns}[T,onlytextwidth]
    \begin{column}{0.46\textwidth}
      \begin{itemize}
        \item Dokumen LPJ pembimbingan dan pendanaan.
        \item Materi dan bank soal yang didistribusikan.
        \item Rekap tryout dan hasil scoring.
        \item Catatan pembahasan serta remedial.
      \end{itemize}
    \end{column}
    \begin{column}{0.48\textwidth}
      \begin{itemize}
        \item Catatan konsultasi akademik.
        \item Rekap strategi belajar dan pilihan jurusan.
        \item Catatan motivasi serta monitoring.
        \item Data aktivitas yang tersedia di StudentPro.
      \end{itemize}
    \end{column}
  \end{columns}
  \vspace{0.22cm}
  \begin{beamercolorbox}[rounded=true,sep=0.19cm,wd=\textwidth]{block body}
    \centering\bfseries Kualitas evaluasi Bab 5 bergantung pada kelengkapan dokumentasi ini.
  \end{beamercolorbox}
\end{frame}

\begin{frame}[shrink=6]{Tidak semua angka operasional telah dicantumkan}
  \kicker{Catatan Keterbatasan Data}
  {\fontsize{10.4}{13}\selectfont
  \renewcommand{\arraystretch}{1.03}
  \begin{tabularx}{\textwidth}{>{\RaggedRight\arraybackslash}p{0.34\textwidth} >{\RaggedRight\arraybackslash}X}
    \toprule
    \textbf{Data yang belum final} & \textbf{Perlakuan dalam laporan} \\
    \midrule
    Jumlah seluruh peserta program & Tidak diperkirakan tanpa rekap final. \\
    Jumlah sesi konsultasi dan tryout & Tidak disamakan dengan jumlah siswa penerima layanan. \\
    Jadwal tanggal rinci kegiatan & Dirangkum setelah arsip kalender diverifikasi. \\
    Capaian penerimaan dan dampak & Dibahas pada Bab 5 menggunakan data hasil. \\
    \bottomrule
  \end{tabularx}}
  \vspace{0.18cm}
  \claim{Keterbatasan dicatat agar laporan tetap jujur dan dapat diperbarui.}
\end{frame}

\begin{frame}{Ringkasan pelaksanaan UTBK 2026}
  \kicker{Apa yang Sudah Berjalan?}
  \begin{columns}[T,onlytextwidth]
    \begin{column}{0.30\textwidth}
      \begin{block}{Belajar}
        Materi, distribusi konten, bank soal, latihan, dan pembahasan.
      \end{block}
    \end{column}
    \begin{column}{0.30\textwidth}
      \begin{block}{Mengukur}
        Diagnostik, tryout berkala, scoring, analisis, dan rekap hasil.
      \end{block}
    \end{column}
    \begin{column}{0.30\textwidth}
      \begin{block}{Mendampingi}
        Konsultasi akademik, strategi jurusan, motivasi, dan monitoring.
      \end{block}
    \end{column}
  \end{columns}
  \vspace{0.30cm}
  \centering{\color{StudentNavy}\bfseries\fontsize{15}{18}\selectfont Pelaksanaan sudah membentuk siklus; Bab 5 menilai efektivitasnya.}
\end{frame}

\begin{frame}{Sumber internal}
  \kicker{Rujukan Bab 4 -- Bagian 1}
  {\color{StudentBlue}\bfseries\fontsize{14.5}{17.5}\selectfont 1. Laporan Pertanggungjawaban Pembimbingan UTBK 2026}\par
  \vspace{0.05cm}
  {\fontsize{11}{13.5}\selectfont Dokumen internal tahun pelajaran 2025/2026. Penyusun: Balya Rochmadi.}\par
  {\color{StudentGray}\fontsize{10}{12.4}\selectfont Sumber bentuk kegiatan, sasaran peserta, konsultasi, monitoring, dan dokumentasi program.}\par
  \vspace{0.34cm}
  {\color{StudentBlue}\bfseries\fontsize{14.5}{17.5}\selectfont 2. Rekap konsultasi akademik UTBK 2026}\par
  \vspace{0.05cm}
  {\fontsize{11}{13.5}\selectfont Catatan diagnostik, strategi belajar, evaluasi tryout, pilihan jurusan/kampus, dan tindak lanjut.}\par
  {\color{StudentGray}\fontsize{10}{12.4}\selectfont Digunakan untuk menyatakan 124 siswa penerima layanan yang terdokumentasi.}\par
\end{frame}

\begin{frame}{Sumber framework dan kesinambungan program}
  \kicker{Rujukan Bab 4 -- Bagian 2}
  {\color{StudentBlue}\bfseries\fontsize{14.5}{17.5}\selectfont 3. Panitia SNPMB}\par
  \vspace{0.05cm}
  {\fontsize{11}{13.5}\selectfont Framework UTBK--SNBT Tahun 2026.}\par
  {\fontsize{10.5}{13}\selectfont\href{https://snpmb.id/fr/}{https://snpmb.id/fr/}}\par
  {\color{StudentGray}\fontsize{10}{12.4}\selectfont Sumber struktur dua kelompok tes dan tujuh subtes yang menjadi acuan pembimbingan.}\par
  \vspace{0.34cm}
  {\color{StudentBlue}\bfseries\fontsize{14.5}{17.5}\selectfont 4. Bab 1--3 Laporan UTBK 2027}\par
  \vspace{0.05cm}
  {\fontsize{11}{13.5}\selectfont Arah program, struktur materi, dan proyeksi penyesuaian menuju UTBK 2027.}\par
  \vfill
  \begin{beamercolorbox}[rounded=true,sep=0.16cm,wd=\textwidth]{block body}
    \centering\bfseries Bab 5 mengevaluasi program SMAN 1 Muntilan berdasarkan bukti pelaksanaan ini.
  \end{beamercolorbox}
\end{frame}

\end{document}
